% --- INICIO DE PLANTILLA PERSONALIZADA ---
\documentclass[oneside,11pt]{Latex/Classes/PhDthesisPSnPDF}

% --- Paquetes Originales ---
\usepackage{blindtext}
\usepackage{actuarialangle} 
\usepackage{tabularx}
\usepackage{float}
\usepackage{multirow, array}
\usepackage[usenames,dvipsnames,svgnames,table]{xcolor}
\usepackage{longtable}
\usepackage{latexsym,amsmath,amssymb,amsfonts}
\usepackage{enumitem}
\usepackage{xparse}
\usepackage{booktabs}
\usepackage[round, sort, numbers]{natbib}  
\usepackage{adjustbox}
\usepackage[utf8]{inputenc}
\usepackage{caption}
\usepackage{graphicx}
\usepackage{hyperref}
\usepackage{cleveref}

% Definición de colores
\definecolor{miblue}{RGB}{68, 114, 196}

% --- CAMBIO CRÍTICO 1: Usar \input para macros en el preámbulo ---
% \include da problemas antes del begin document. \input es seguro.
\input{Latex/Macros/MacroFile1}

% --- Definiciones de seguridad (por si la clase falla al cargar alguna variable) ---
\providecommand{\facultad}[1]{\def\lafacultad{#1}}
\providecommand{\escudofacultad}[1]{\def\elescudo{#1}}
\providecommand{\degree}[1]{\def\elgrado{#1}}
\providecommand{\director}[1]{\def\eldirector{#1}}
\providecommand{\lugar}[1]{\def\ellugar{#1}}
\providecommand{\degreedate}[1]{\def\lafecha{#1}}

% --- Configuración para bloques de código de R (Pandoc) ---
\usepackage{color}
\usepackage{fancyvrb}
\newcommand{\VerbBar}{|}
\newcommand{\VERB}{\Verb[commandchars=\\\{\}]}
\DefineVerbatimEnvironment{Highlighting}{Verbatim}{commandchars=\\\{\}}
% Add ',fontsize=\small' for more characters per line
\usepackage{framed}
\definecolor{shadecolor}{RGB}{248,248,248}
\newenvironment{Shaded}{\begin{snugshade}}{\end{snugshade}}
\newcommand{\AlertTok}[1]{\textcolor[rgb]{0.94,0.16,0.16}{#1}}
\newcommand{\AnnotationTok}[1]{\textcolor[rgb]{0.56,0.35,0.01}{\textbf{\textit{#1}}}}
\newcommand{\AttributeTok}[1]{\textcolor[rgb]{0.77,0.63,0.00}{#1}}
\newcommand{\BaseNTok}[1]{\textcolor[rgb]{0.00,0.00,0.81}{#1}}
\newcommand{\BuiltInTok}[1]{#1}
\newcommand{\CharTok}[1]{\textcolor[rgb]{0.31,0.60,0.02}{#1}}
\newcommand{\CommentTok}[1]{\textcolor[rgb]{0.56,0.35,0.01}{\textit{#1}}}
\newcommand{\CommentVarTok}[1]{\textcolor[rgb]{0.56,0.35,0.01}{\textbf{\textit{#1}}}}
\newcommand{\ConstantTok}[1]{\textcolor[rgb]{0.00,0.00,0.00}{#1}}
\newcommand{\ControlFlowTok}[1]{\textcolor[rgb]{0.13,0.29,0.53}{\textbf{#1}}}
\newcommand{\DataTypeTok}[1]{\textcolor[rgb]{0.13,0.29,0.53}{#1}}
\newcommand{\DecValTok}[1]{\textcolor[rgb]{0.00,0.00,0.81}{#1}}
\newcommand{\DocumentationTok}[1]{\textcolor[rgb]{0.56,0.35,0.01}{\textbf{\textit{#1}}}}
\newcommand{\ErrorTok}[1]{\textcolor[rgb]{0.64,0.00,0.00}{\textbf{#1}}}
\newcommand{\ExtensionTok}[1]{#1}
\newcommand{\FloatTok}[1]{\textcolor[rgb]{0.00,0.00,0.81}{#1}}
\newcommand{\FunctionTok}[1]{\textcolor[rgb]{0.00,0.00,0.00}{#1}}
\newcommand{\ImportTok}[1]{#1}
\newcommand{\InformationTok}[1]{\textcolor[rgb]{0.56,0.35,0.01}{\textbf{\textit{#1}}}}
\newcommand{\KeywordTok}[1]{\textcolor[rgb]{0.13,0.29,0.53}{\textbf{#1}}}
\newcommand{\NormalTok}[1]{#1}
\newcommand{\OperatorTok}[1]{\textcolor[rgb]{0.81,0.36,0.00}{\textbf{#1}}}
\newcommand{\OtherTok}[1]{\textcolor[rgb]{0.56,0.35,0.01}{#1}}
\newcommand{\PreprocessorTok}[1]{\textcolor[rgb]{0.56,0.35,0.01}{\textit{#1}}}
\newcommand{\RegionMarkerTok}[1]{#1}
\newcommand{\SpecialCharTok}[1]{\textcolor[rgb]{0.00,0.00,0.00}{#1}}
\newcommand{\SpecialStringTok}[1]{\textcolor[rgb]{0.31,0.60,0.02}{#1}}
\newcommand{\StringTok}[1]{\textcolor[rgb]{0.31,0.60,0.02}{#1}}
\newcommand{\VariableTok}[1]{\textcolor[rgb]{0.00,0.00,0.00}{#1}}
\newcommand{\VerbatimStringTok}[1]{\textcolor[rgb]{0.31,0.60,0.02}{#1}}
\newcommand{\WarningTok}[1]{\textcolor[rgb]{0.56,0.35,0.01}{\textbf{\textit{#1}}}}

%%%%%%%%%%%%%%%%%%%%%%%%%%%%%%%%%%%%%%%%%%%%%%%%%%%%%%%%%%%%%%%%%%%%%%%%%%%%%%%%
%                         DATOS (Conectados a RMarkdown)                       %
%%%%%%%%%%%%%%%%%%%%%%%%%%%%%%%%%%%%%%%%%%%%%%%%%%%%%%%%%%%%%%%%%%%%%%%%%%%%%%%%
\title{El Consumo de alcohol per capital y los datos que no se miden}
\author{Andrea Figuera \\ Eduardo Espinoza}

% Lógica para Facultad: Si la defines en el YAML la usa, si no, usa la default

\facultad{Facultad de Ciencias Económicas y Sociales \\ Escuela de
Estadistica y Ciencias Actuariales}
  \escudofacultad{Latex/Classes/Escudos/faces}

\degree{Computación I}
\director{Jesus Ochoa \\ Oliver Triveño}
\degreedate{Abril 2026} 
\lugar{Caracas}

\portadatrue 

% Metadatos PDF
\keywords{}
\subject{}

% --- CORRECCIÓN PARA PANDOC NUEVO (Imágenes) ---
\makeatletter
\providecommand{\pandocbounded}[1]{#1}
\makeatother


%%%%%%%%%%%%%%%%%%%%%%%%%%%%%%%%%%%%%%%%%%%%%%%%%%%%%
%                   DOCUMENTO                       %
%%%%%%%%%%%%%%%%%%%%%%%%%%%%%%%%%%%%%%%%%%%%%%%%%%%%%
\begin{document}

\maketitle

%%%%%%%%%%%%%%%%%%%%%%%%%%%%%%%%%%%%%%%%%%%%%%%%%%%%%
%                  PRÓLOGO                          %
%%%%%%%%%%%%%%%%%%%%%%%%%%%%%%%%%%%%%%%%%%%%%%%%%%%%%
\frontmatter

% Puedes agregar dedicatorias desde el YAML si quieres, o dejarlas fijas aquí

%%%%%%%%%%%%%%%%%%%%%%%%%%%%%%%%%%%%%%%%%%%%%%%%%%%%%
%                   ÍNDICES                         %
%%%%%%%%%%%%%%%%%%%%%%%%%%%%%%%%%%%%%%%%%%%%%%%%%%%%%
\setcounter{secnumdepth}{3}
\setcounter{tocdepth}{3}

\tableofcontents



\cleardoublepage

  \cleardoublepage       % Empieza en página derecha (estándar en tesis)
  \chapter*{Resumen}     % Crea el título "Resumen" sin numeración (Capítulo *)
  \addcontentsline{toc}{chapter}{Resumen} % (Opcional) Lo añade al índice
  
  Resumen del trabajo de
investigación\ldots{}             % Aquí se pega el texto que escribiste en el YAML
  
  \cleardoublepage

%%%%%%%%%%%%%%%%%%%%%%%%%%%%%%%%%%%%%%%%%%%%%%%%%%%%%
%                   CONTENIDO (El Corazón)          %
%%%%%%%%%%%%%%%%%%%%%%%%%%%%%%%%%%%%%%%%%%%%%%%%%%%%%
\mainmatter
\def\baselinestretch{1.5}

% --- CAMBIO CRÍTICO 2: Aquí RMarkdown inyecta todo tu texto ---
\chapter*{Introducción}

Introducción del trabajo de investigación

\chapter{Planteamiento del Problema}

Cuando se habla del consumo de alcohol al nivel mundial, los resultados
de un análisis generalmente son presentados según qué tan por encima o
por debajo esté el promedio de cierta población estudiada en relación
con el \textbf{indicador transversal} utilizado por la Organización
Mundial de la Salud (OMS). El consumo de alcohol per cápital es un tema
en el que se evalúa el riesgo y el punto donde el consumo se categoriza
como un tema crítico para la salud pública, al momento de determinarlo
se choca directamento con el \textbf{reduccionismo estadístico}, el cual
hace referencia al como la interpretación de los datos viene presentada
bajo \textbf{indicadores} estadísticos que buscan reducir la realidad
objetiva de poblaciones enteras a un simple promedio.

Entes oficiales de salud pública como la OMS o la Organización
Panamericana de la Salud (OPS) usan herramientas estadísticas para la
recolección, interpretación y presentación de los datos; la OMS
establece una medida estandar para el análisis global del consumo de
alcohol que define de la siguiente manera:

\begin{quote}
El consumo total de alcohol per cápita (APC) se define como la cantidad
total (suma del promedio de tres añoS registrado y el promedio de tres
años de APC no registrado, ajustado por el consumo turístico promedio de
tres años) de alcohol consumido por adulto (15+ años) durante un año
calendario, en litros de alcohol puro. (OMS, 2018, p.~411)
\end{quote}

Esta definición revela la complejidad que puede esconderse detrás de un
solo numero. Al involucrar tantas variables el APC pasa de ser solo un
promedio a convertirse mas en una estimación sujeta a un gran margen de
error. Sin embargo, en el discurso público, esta cifra suele presentarse
como una verdad absoluta, ignorando que el promedio es una medida
extremadamente sensible a valores atípicos, extremos y faltantes. Un
promedio suele ser una medida muy débil al momento de trabajar con datos
que contienen mayor dispersión.

Al hablar del consumo de alcohol generalmente se encuentra un sesgo muy
grande entre los datos: si un grupo pequeño consume de forma masiva, el
promedio subirá, estigmatizando a toda una población por el
comportamiento de una minoría, alejandose de la realidad o peor aún, se
asume que el riesgo está distribuido de forma equitativa cuando en
realidad está concentrado en sectores específicos que el promedio no
logra mostrar. En ciertos casos la mediana suele ser más representativa
pero los entes políticos prefieren el promedio por ser este más
manipulable. Si un país no registra bien sus datos, el promedio
simplemente ignora esos vacíos, esto genera un \textbf{sesgo de
selección}, el promedio parece saludable solo porque no se esta contando
los datos que no se puden medir. Entonces esto dice que ``lo que no se
mide, no existe''. Siendo los valores faltantes un peligro. Muchas
veces, los países con las crisis de consumo más graves tienen los
sistemas de registro más deficientes. Cuando un reporte ignora estos
huecos de información, el promedio nacional se ve ``limpio'',
permitiendo que se construyan narrativas basadas en un sesgo de
selección. Al final, exponer estos datos sin explicar la variabilidad
técnica (como el ancho de los intervalos de confianza) convierte la
estadística en una herramienta que oculta las verdaderas crisis de salud
detrás de una falsa sensación de precisión.

A todo lo antes expuesto le queda la frase popular conocida en la
estadística; ``No cruces un río que tiene un promedio de un metro de
profundidad, porque podrías ahogarte en la parte que mide tres metros''

El promedio es la medida que los gobiernos suelen usar para decir si un
país tiene un elevado consumo de alcohol en comparación con el promedio
mundial de la OMS. También se utiliza la \textbf{Prevalencia de
Consumo}, que mide qué porcentaje de la población bebió en el último
año, y un indicador más técnico llamado \textbf{DALY}, que estima
cuántos años de vida saludable pierde una comunidad por problemas con la
bebida. El detalle está cuando estas cifras se presentan de forma
distinta a lo que son, \textbf{estimaciones}. Al enfocarse solo en estos
números finales, se deja de lado la incertidumbre y el error que hay
detrás de cada dato, permitiendo que se construyan discursos donde un
promedio bajo oculte una crisis de salud real en sectores específicos.

Como dicta la conocida paradoja estadística: si una persona se come un
pollo entero y otra no come nada, el promedio indicará que cada una
consumió medio pollo. En la práctica, este sesgo puede ser utilizado
políticamente para proyectar una falsa sensación de éxito o estabilidad,
permitiendo que un gobierno declare un bajo consumo en su región
mientras ignora deliberadamente a los sectores con consumos explosivos o
problemáticos que quedan ocultos tras la cifra central.

Esta ``mentira estadística'' sustenta los factores críticos que afectan
la veracidad de los datos recolectados:

Un promedio carece de valor científico si no se analiza su margen de
error.

Una región puede reportar un consumo bajo, pero si el ancho de su
intervalo de confianza es excesivamente grande, esto no refleja un
comportamiento social claro, sino una deficiencia en la calidad de la
recolección de los datos o una variabilidad inmanejable.

El error más común en la comparación internacional es asumir que
distintos promedios implican distintas realidades. Si al comparar dos
naciones sus intervalos de confianza se solapan significativamente,
estadísticamente no existe evidencia suficiente para afirmar que una
consume más que la otra. El solapamiento es la prueba de que el
promedio, por sí solo, es insuficiente para establecer jerarquías de
consumo.

La ausencia de registros en ciertos periodos o sectores poblacionales
desplaza el promedio hacia valores que no representan la realidad. El
uso de bases de datos con vacíos informativos permite construir
narrativas que benefician agendas institucionales en lugar de reflejar
la salud pública real.

Un caso real que ilustra esta problemática se observa en los reportes
globales de la Organización Mundial de la Salud (OMS). En regiones con
sistemas de vigilancia débiles, se reportan promedios nacionales que
subestiman el riesgo real debido a que la incertidumbre de la medición
es tan alta que el dato central deja de ser representativo. De igual
forma, comparaciones clásicas entre Europa del Este y Europa Occidental
a menudo revelan que, aunque los titulares de prensa resalten
diferencias en los promedios, los intervalos de confianza de estos
países suelen solaparse, indicando que las diferencias podrían no ser
estadísticamente significativas.

Esto responde a que la variabilidad del consumo no es constante ni
uniforme entre las naciones.Por lo tanto, este trabajo de investigación
se propone cuestionar la validez de las comparaciones tradicionales de
consumo de alcohol. El problema no radica solo en la cantidad consumida,
sino en la incertidumbre estadística que se desprende de los datos. El
desafío es, entonces la implementación de algoritmos que permitan una
gestión rigurosa de los datos faltantes \textbf{(missing data)},
entendiendo que su ausencia no es aleatoria, sino un indicador de la
debilidad de los sistemas de vigilancia. Más allá de la implementación
de medidas de tendencia central, este trabajo se enfoca en la
visualización de densidades de probabilidad y el procesamiento de
intervalos de confianza. El objetivo es demostrar mediante el análisis
de datos que, ante una alta varianza y un solapamiento significativo de
las distribuciones, el promedio pierde su capacidad explicativa.
Actuando en colaboración con la ciencia computacional: se requiere
procesar y visualizar estos niveles de solapamiento y dispersión para
demostrar que, sin un análisis riguroso de la incertidumbre, la
estadística deja de ser una herramienta de precisión para convertirse en
una narrativa que oculta las verdaderas crisis de consumo.\\

Argumento final. Debido a esto, se plantea las siguientes preguntas de
investigación:\\

\textbf{1.} ¿Cómo influyen los valores atípicos y los datos faltantes en
la representatividad del promedio del consumo de alcohol?\\

\textbf{2.} ¿Qué tan confiables son los promedios globales si
consideramos que los márgenes de error varían drásticamente entre
regiones, y cómo afecta esta diferencia la precisión de lo que se
reporta?\\

\textbf{3.} ¿En qué medida el solapamiento de los intervalos de
confianza invalida las comparaciones y jerarquías de consumo
tradicionales?\\

\section{Justificación}

La presente investigación se fundamenta en la necesidad de pasar de una
estadística descriptiva básica hacia un análisis de datos crítico que
integre la \textbf{incertidumbre} como una variable fundamental para ser
presentada a la sociedad. En el ámbito académico de la \textbf{Escuela
de Estadística y Ciencias Actuariales (EECA)}, es necesario demostrar
que el rol del analista no es simplemente procesar cifras, sino
cuestionar la representatividad de la información que sustenta las
políticas públicas globales. La importancia de este análisis radica en
desmontar el uso político de las medidas de tendencia central. Un
gobierno puede declarar una gestión exitosa basándose en la reducción de
un promedio, ignorando deliberadamente que dicha cifra puede ser el
resultado de datos altamente sesgados o con una incertidumbre tan
elevada que la comparación pierde validez científica. Al estudiar el
solapamiento de los intervalos, esta investigación busca devolverle la
rigurosidad a la interpretación de los datos, demostrando que la
estadística no debe servir para simplificar la realidad, sino para
revelar su complejidad.

\textbf{Relevancia Técnica y Delimitación del Alcance} Desde una
perspectiva teórica, el trabajo contribuye al campo de la estadística
social al demostrar que el promedio, como medida de tendencia central
única, es insuficiente para caracterizar el consumo de alcohol. Al
introducir conceptos como el \textbf{solapamiento de intervalos de
confianza y la variabilidad}, esta investigación enriquece el debate
académico sobre cómo deben compararse las realidades de distintas
naciones, evitando conclusiones apresuradas que ignoran la dispersión de
los datos.

Desde el punto de vista computacional, este proyecto se justifica al
utilizar herramientas de programación (R y Python) para extraer y tratar
la \textbf{incertidumbre} declarada en los reportes oficiales. Dado que
la investigación se desarrolla en una etapa inicial de la formación
académica, el enfoque no pretende realizar estimaciones propias ni
modelos de regresión, sino centrarse exclusivamente en el
\textbf{procesamiento y visualización del solapamiento de los intervalos
de confianza} ya proporcionados en el dataset. Esta delimitación es
clave: el valor agregado de este trabajo radica en la capacidad de
convertir tablas de datos crudos en \textbf{visualizaciones de
densidades de probabilidad} que evidencien visualmente cuándo una
diferencia entre países es estadísticamente significativa y cuándo es
simplemente producto de la variabilidad del indicador transversal.

\textbf{Relevancia Social y Ética} Socialmente, el trabajo se justifica
al exponer cómo el reduccionismo estadístico puede ser utilizado para
construir narrativas que ocultan crisis sanitarias reales. Al procesar
los intervalos de confianza y los datos faltantes presentes en el
dataset, se busca demostrar que un promedio nacional ``limpio'' puede
ser el resultado de sistemas de registro deficientes, permitiendo que se
ignore deliberadamente a los sectores con consumos problemáticos.
Identificar vulnerabilidades ocultas, a la hora de tomar decisiones y
entender el ancho de la incertidumbre permite reconocer en qué regiones
la recolección de datos es deficiente o dónde el consumo es tan dispar
que un promedio bajo no garantiza la ausencia de riesgos. Esto permite
transitar de una política basada en cifras generales a una basada en la
realidad de los intervalos, protegiendo así a los sectores poblacionales
que quedan invisibilizados en los reportes tradicionales.

Así la presente investigación se plantea por la necesidad imperativa de
dotar a la salud pública de herramientas de análisis que trasciendan la
interpretación superficial de los datos. En un contexto global donde la
información estadística es utilizada con frecuencia para validar
narrativas políticas, este estudio aporta un marco crítico que devuelve
la importancia a la \textbf{variabilidad e incertidumbre} como
componentes esenciales de la verdad científica.

\textbf{Aporte Académico} Finalmente, la \textbf{justificación
metodológica} se sustenta en el uso de herramientas computacionales de
vanguardia. La complejidad de procesar un conjunto de datos extenso, que
abarca múltiples años y regiones con valores faltantes, hace que el uso
de \textbf{Python o R} no sea opcional, sino un requisito técnico. El
manejo de estructuras de datos y la generación de visualizaciones
avanzadas permiten exponer visualmente el solapamiento que los métodos
manuales no logran captar. Así, este proyecto no solo cumple con una
exigencia académica, sino que demuestra cómo la computación aplicada a
la estadística es el único camino para desarticular las ``mentiras'' que
surgen de un análisis de datos incompleto.

Este proyecto contribuye al repositorio \textbf{``Análisis de consumo de
alcohol e incertidumbre (Computación 1 - UCV)''}. Aporta un modelo de
análisis que no solo lee cifras, sino que evalúa la veracidad de las
comparaciones internacionales mediante el rigor computacional,
rescatando la hipótesis de que la estadística debe servir para revelar
realidades, no para ocultarlas tras una falsa sensación de precisión.\\

\section{Objetivos}

\subsection{Objetivo General}
\begin{itemize}
  \item{Cuestionar la validez de las comparaciones tradicionales del consumo de alcohol a nivel mundial mediante el procesamiento computacional de los intervalos de confianza proporcionados en el dataset, con el fin de evidenciar el solapamiento de datos y la incertidumbre estadística que el promedio suele invisibilizar.}
\end{itemize}

\subsection{Objetivos Específicos}
\begin{itemize}
  \item{Extraer y depurar los registros de consumo de alcohol y sus intervalos de confianza, gestionando los valores faltantes y contrastando la media frente a la mediana para identificar el sesgo generado por valores atípicos (outliers) en el promedio.}
  \item{Evaluar la variabilidad en la precisión de las estimaciones analizando el ancho de los márgenes de error para cuestionar la consistencia de los reportes globales}
  \item{Identificar y visualizar los niveles de solapamiento de los intervalos de confianza, con el fin de determinar regiones donde las diferencias en los promedios carecen de verdad estadística, culminando en la generación de un reporte técnico reproducible que evidencie la incertidumbre de los datos.}
\end{itemize}

\section{Cobertura}

\subsection{Horizontal}

La cobertura horizontal del presente estudio abarca las características
intrínsecas del consumo de alcohol registradas en la base de datos
alcohol\_data.csv. Se consideran las \textbf{dimensiones descriptivas}
del registro (país (country), código iso3, y año
(year)).\textbf{Adicionalmente, se incorporará una variable derivada
para agrupar los países en regiones geográficas, facilitando un análisis
comparativo a mayor escala}. Se incluyen las \textbf{métricas de
tendencia central}, específicamente el consumo de alcohol per cápita en
litros (alcohol\_liters\_per\_capita).Finalmente, se integran las
\textbf{variables de incertidumbre técnica}, analizando el límite
inferior (lower\_ci), el límite superior (upper\_ci) y el ancho del
intervalo de confianza (ci\_width) para cuestionar la precisión de los
reportes globales.\\

\begin{table}[ht]
\centering
\caption{Variables consideradas en el análisis de consumo e incertidumbre}
\label{tab:variables}
% Ajustamos el ancho a 12cm para que quepa bien el texto
\resizebox{15cm}{!} {
% Definimos 3 columnas centradas (c c c)
\begin{tabular}{ccc}
  \hline
  \textbf{NOMBRE DE LA VARIABLE} & \textbf{TIPO DE VARIABLE} & \textbf{DESCRIPCIÓN} \\
  \hline
  país (`country`) & Categórica (Dimensión) & Nombre del país reportado. \\
  año (`year`) & Numérica Entera (Tiempo) & Año del reporte (2000-2019). \\
  consumo (`alcohol\_liters\_per\_capita`) & Numérica Continua (Métrica) & Consumo per cápita en litros. \\
  ancho del intervalo (`ci\_width`) & Numérica Continua (Incertidumbre) & Medida de la precisión del reporte. \\
  región (derivada) & Categórica (Agrupación) & Región geográfica (creada para análisis). \\
  \hline
\end{tabular}
}
\end{table}

\subsection{Vertical}

La cobertura vertical del trabajo de investigación comprende un total de
4324 observaciones (registros). Cada fila representa un dato anual de
consumo por país. Para efectos de esta investigación, se realizó una
delimitación poblacional filtrando únicamente los registros
correspondientes a \textbf{``Ambos sexos''} (Both sexes). El alcance
temporal de los datos es longitudinal, abarcando desde el año
\textbf{2000 hasta el año 2019}. Esta delimitación permite un análisis
comparativo y \textbf{evolutivo} para observar las variaciones en los
reportes de salud pública a lo largo de estas dos décadas.\\

\section{Periodo de Referencia}

Para la presente investigación, el periodo de referencia efectivo para
el análisis descriptivo y de incertidumbre técnica se estableció en el
intervalo \textbf{2000-2019}. Este rango de dos décadas fue seleccionado
estratégicamente debido a que presenta la mayor \textbf{consistencia y
precisión} en los reportes de salud global sobre consumo de alcohol per
cápita. A partir del año 2000, los métodos de recopilación de datos y
los modelos de estimación de intervalos de confianza fueron
estandarizados, lo que reduce los márgenes de error (\texttt{ci\_width})
y permite una comparación longitudinal más robusta entre países y
regiones.El procesamiento de la información se realizó con corte a Abril
de 2026\\

\section{Variables}

\textbf{La variable principal en estudio (Variable Dependiente/Objetivo)
es el consumo
(\texttt{alcohol\textbackslash{}\_liters\textbackslash{}\_per\textbackslash{}\_capita})},
que representa el promedio de litros de alcohol puro consumidos por
persona al año en la población de Ambos sexos (\texttt{Both\ sexes}).
Esta métrica de tendencia central es fundamental para evaluar las
tendencias globales de salud pública.

\textbf{Como variables explicativas o independientes se analizan:}

\begin{enumerate}
\def\labelenumi{\arabic{enumi}.}
\tightlist
\item
  \textbf{Variables Temporales y Geográficas:}

  \begin{itemize}
  \tightlist
  \item
    \textbf{Año (\texttt{year})}: Año del reporte (2000-2019).
  \item
    \textbf{País (\texttt{country})}: Nombre del país reportado.
  \item
    \textbf{Código ISO3 (\texttt{iso3})}: Código estandarizado de tres
    letras del país.
  \end{itemize}
\item
  \textbf{Variable Geográfica Derivada:}

  \begin{itemize}
  \tightlist
  \item
    \textbf{Región (derivada)}: Agrupación geográfica de países (creada
    para este análisis, ej: ``Américas'', ``Europa'', ``África''). Esta
    variable permite un análisis comparativo y de consistencia a una
    escala mayor, facilitando la identificación de patrones regionales
    de incertidumbre técnica.
  \end{itemize}
\item
  \textbf{Variables de Incertidumbre Técnica:}

  \begin{itemize}
  \tightlist
  \item
    \textbf{Intervalo de Confianza (Lower CI / Upper CI)}: Límites
    inferior (\texttt{lower\textbackslash{}\_ci}) y superior
    (\texttt{upper\textbackslash{}\_ci}) que delimitan el rango de la
    estimación.
  \item
    \textbf{Ancho del Intervalo (\texttt{ci\textbackslash{}\_width})}:
    Medida de la precisión técnica, calculada como la diferencia entre
    los límites superior e inferior. Un ancho mayor indica una menor
    consistencia en el reporte global./
  \end{itemize}
\end{enumerate}

\chapter{Marco Teórico}

\section{Bases Teóricas}

El presente capítulo establece los fundamentos estadísticos y
actuariales necesarios para el análisis del consumo de alcohol a nivel
global y la cuantificación de su incertidumbre. Todo el marco conceptual
se ajusta estrictamente a las variables observadas en la base de datos
objeto de este estudio.

\subsection{Consumo de Alcohol Total Per Cápita (APC)}

El indicador principal de este estudio es el Consumo de Alcohol Total
Per Cápita (APC), representado en el conjunto de datos por la variable
\texttt{alcohol\_liters\_per\_capita}. Técnicamente, se define como la
cantidad total (en litros) de alcohol puro consumido por la población en
un país durante un año específico.

Desde la perspectiva estadístico-actuarial, el valor reportado como APC
es un \textbf{estimador puntual} de la media poblacional (\(\mu\)). Dado
que es imposible conocer el consumo exacto y absoluto de cada individuo
en una nación (debido a la existencia de consumo no registrado, como la
producción artesanal o el mercado ilícito), este valor se construye a
partir de modelos de estimación que combinan fuentes determinísticas y
estocásticas.

\subsection{Incertidumbre y Componente Estocástico}

Debido a que el APC incluye datos que no se miden de forma directa y
oficial, no puede considerarse un valor determinístico perfecto. Por
ello, el análisis requiere el uso de
\textbf{Intervalos de Confianza (IC)} para cuantificar el riesgo de
error en la medición.

Para cada estimación del APC (\(\hat{\theta}\)), la base de datos
proporciona un límite inferior (\texttt{lower\_ci}) y un límite superior
(\texttt{upper\_ci}), definiendo el intervalo para cada observación:

\[IC = [\hat{\theta}_{lower}, \hat{\theta}_{upper}]\]

La diferencia entre el límite superior y el inferior se define como la
\textbf{Amplitud del Intervalo} (representada por la variable
\texttt{ci\_width}):

\[ci\_width = \hat{\theta}_{upper} - \hat{\theta}_{lower}\]

Esta amplitud es la medida directa de la
\textbf{incertidumbre estocástica}. Un valor elevado de
\texttt{ci\_width} indica una mayor variabilidad y menor precisión en
los datos reportados por ese país, lo cual constituye el núcleo
analítico de la presente investigación.

\subsection{Medidas de Tendencia Central y Asimetría}

Para evaluar el comportamiento del consumo a nivel agregado, es
fundamental analizar la forma de su distribución. Para ello, se
contrastará la \textbf{Media Muestral} (\(\bar{x}\)) con la
\textbf{Mediana} (\(\tilde{x}\)) de la variable
\texttt{alcohol\_liters\_per\_capita}.

\begin{itemize}
\tightlist
\item
  \textbf{Media Muestral:} Representa el centro de gravedad o equilibrio
  de la distribución de los datos. Matemáticamente, se define como la
  suma de todos los valores observados dividida entre el tamaño de la
  muestra (n). Es una medida altamente sensible a valores extremos.
  \[\bar{x} = \frac{1}{n} \sum_{i=1}^{n} x_i\]
\item
  \textbf{Mediana:} Es el valor posicional que divide la serie de datos
  ordenados exactamente en dos partes iguales (50\% superior y 50\%
  inferior). A diferencia de la media, es una medida robusta, lo que
  significa que no se ve afectada por valores atípicos (outliers) de
  consumo extremo.
\end{itemize}

La relación entre ambas medidas permite identificar sesgos
distribucionales. Si \(\bar{x} > \tilde{x}\), existe una asimetría
positiva, lo que indicaría que un grupo reducido de países con consumos
extremos eleva el promedio general, alejándolo del comportamiento típico
de la mayoría de las naciones.

\subsection{Criterio de Solapamiento de Intervalos}

Para determinar estocásticamente si el consumo entre dos países, o entre
dos años distintos, presenta diferencias estadísticamente
significativas, se empleará el análisis del solapamiento de los
intervalos de confianza (\texttt{lower\_ci} y \texttt{upper\_ci}).

Se establecen los siguientes criterios:

\begin{itemize}
    \item Existe una \textbf{diferencia significativa} si el límite superior de la estimación de un grupo $A$ es estrictamente menor al límite inferior de la estimación de un grupo $B$ ($Upper\_CI_A < Lower\_CI_B$).
    \item Existe \textbf{solapamiento} si los intervalos comparten un rango de valores. En este escenario, la diferencia observada en los estimadores puntuales (el APC) podría deberse puramente a la variabilidad aleatoria y a la falta de precisión inherente a los datos que no se miden, impidiendo concluir con certeza que una población consume más que otra.
\end{itemize}

\subsection{Robustez Estadística}

En el análisis de datos sobre consumo de alcohol, la \textbf{robustez}
se define como la propiedad de un estadístico de permanecer inalterado o
mínimamente afectado ante la presencia de valores atípicos
(\textit{outliers}) o errores de medida en la muestra.

En esta investigación, la robustez es el criterio principal para
preferir la \textbf{Mediana} sobre la \textbf{Media Muestral} al momento
de describir el comportamiento típico de una región. Mientras que la
media se desplaza hacia los extremos cuando existen países con consumos
desproporcionadamente altos, la mediana mantiene su posición central,
proporcionando una visión más estable del fenómeno ante la incertidumbre
de los datos.

\subsection{Sesgo de Información y Variabilidad}

El sesgo en este contexto se refiere a la distorsión de la medida
central respecto a la realidad de la mayoría de la población. Dado que
el dataset incluye el componente de incertidumbre (\texttt{ci\_width}),
es necesario definir dos conceptos clave para el análisis de resultados:

\begin{enumerate}
    \item \textbf{Sesgo por Valores Extremos:} Se identifica mediante la diferencia entre la media y la mediana ($Sesgo = \bar{x} - \tilde{x}$). Un valor positivo indica que el promedio está "inflado" por países con registros inusualmente altos.
    \item \textbf{Variabilidad de la Estimación:} Es el grado de dispersión representado por el ancho del intervalo de confianza. A mayor variabilidad, menor es la confiabilidad de la media como indicador único de riesgo para el diseño de políticas públicas.
\end{enumerate}

\hfill\break

\chapter{Marco Metódico}

En esta sección se presentan las metodologías, procedimientos
algorítmicos y pruebas estadísticas implementadas para validar las
hipótesis planteadas en el capítulo anterior. Se detalla el tratamiento
de la fuente de datos \texttt{alcohol\_data.csv}, las técnicas de
limpieza de datos (Data Wrangling) y los métodos de comparación de
intervalos utilizados para determinar la significancia estadística de
las variaciones en el consumo.

Un aspecto fundamental de este diseño es la estratificación de los datos
por regiones geográficas, lo que permite identificar patrones de consumo
y niveles de incertidumbre diferenciados. Asimismo, con el objetivo de
garantizar la estabilidad de las comparaciones y evitar sesgos
coyunturales derivados de la crisis sanitaria global, la investigación
toma como ventana de observación los últimos cinco años previos al
fenómeno de la pandemia (2015-2019). Esta delimitación temporal permite
analizar el comportamiento del APC en un contexto de normalidad
socioeconómica, asegurando que la incertidumbre detectada responda a
deficiencias estructurales en los registros y no a las alteraciones
atípicas en el consumo y la vigilancia epidemiológica ocurridas durante
el periodo del COVID-19.

El análisis se ha realizado utilizando el lenguaje de programación R
bajo el entorno de desarrollo RStudio, empleando el paradigma
``Tidyverse'' para la manipulación eficiente de datos y la librería
``ggplot2'' para la visualización de los intervalos de confianza,
facilitando así el estudio del comportamiento estocástico del consumo de
alcohol a nivel regional y global.De igual manera se utilizara el
lenguaje Python, mediante la librería Streamlit, para el desarrollo de
un dashboard interactivo. Esta herramienta permite la visualización
dinámica de los intervalos de confianza y el consumo por regiones,
facilitando así el estudio del comportamiento de los datos. esto es lo
que tengo el marco metodico y luego debo poner las bases metodicas

\section{Bases metódicas utilizadas en el trabajo de investigación}

La presente investigación es de tipo descriptiva. Su objetivo central es
caracterizar el fenómeno del consumo de alcohol per cápita y,
fundamentalmente, describir la evolución de la incertidumbre técnica
asociada a estos reportes globales entre los años 2015 y 2019. No se
busca establecer relaciones de causa y efecto, sino realizar una
``radiografía'' precisa de la consistencia de los datos a través del
análisis de medidas de tendencia central (media) y dispersión
(desviación estándar, ancho del intervalo de confianza
\texttt{ci\textbackslash{}\_width}).

El diseño de la investigación es longitudinal, específicamente un
análisis de datos de panel. Esto se debe a que se analizan las mismas
variables (consumo e incertidumbre) para las mismas unidades de
observación (países) a lo largo de un periodo de tiempo continuo
(2015-2019). Este diseño permite observar no solo el estado actual, sino
los cambios y la evolución temporal de la precisión en los reportes de
salud pública.

Dado que la investigación se basa exclusivamente en el análisis de una
base de datos preexistente (\texttt{alcohol\textbackslash{}\_data.csv}),
se clasifica como una investigación documental o basada en fuentes
secundarias. No hubo recolección directa de datos primarios en campo; el
trabajo consistió en el procesamiento, limpieza y análisis estadístico
de información ya recopilada y validada por una autoridad global.

En esta sección se detallan los procedimientos algorítmicos y las
técnicas estadísticas que constituyen el sustento operativo de la
investigación, ejecutados secuencialmente para alcanzar los objetivos
planteados\\

\subsection{Procedimiento de Data Wrangling (Procesamiento de Datos)}

El tratamiento de la fuente de datos
\texttt{alcohol\textbackslash{}\_data.csv} se realizó siguiendo el
paradigma ``Tidyverse'' en R, estructurado en las siguientes fases
operativas:

\begin{enumerate}
\def\labelenumi{\arabic{enumi}.}
\tightlist
\item
  \textbf{Carga y Validación de Estructura:} Importación del conjunto de
  datos y verificación de tipos de variables (categóricas y numéricas
  continua/entera) mediante las funciones \texttt{read.csv()} y
  \texttt{str()}.
\item
  \textbf{Delimitación Poblacional y Temporal:} Implementación de
  filtros algorítmicos con la función \texttt{filter()} de la librería
  \texttt{dplyr} para seleccionar exclusivamente los registros de ``Both
  sexes'' (Ambos sexos) y acotar el análisis a la ventana de observación
  efectiva (2015-2019).
\item
  \textbf{Selección de Atributos:} Uso de \texttt{select()} para
  conservar únicamente las variables críticas para el estudio:
  \texttt{country}, \texttt{iso3}, \texttt{year},
  \texttt{alcohol\textbackslash{}\_liters\textbackslash{}\_per\textbackslash{}\_capita},
  \texttt{lower\textbackslash{}\_ci}, \texttt{upper\textbackslash{}\_ci}
  y \texttt{ci\textbackslash{}\_width}.
\item
  \textbf{Estratificación Geográfica (Creación de Variable Derivada):}
  Desarrollo de un algoritmo de recodificación (mediante
  \texttt{mutate()} y \texttt{case\_when()}) para asignar a cada país su
  respectiva región geográfica (ej: ``África'', ``Américas'', ``Sureste
  Asiático'', ``Europa'', ``Mediterráneo Oriental'', ``Pacífico
  Occidental''), conveniente para el estudio.\\
\end{enumerate}

\begin{Shaded}
\begin{Highlighting}[]
\FunctionTok{library}\NormalTok{(tidyverse)}
\FunctionTok{library}\NormalTok{(knitr)}

\NormalTok{raw\_data\_alcohol }\OtherTok{\textless{}{-}} \FunctionTok{read.csv}\NormalTok{(}\StringTok{"alcohol\_data.csv"}\NormalTok{) }\CommentTok{\#Importación del conjunto de datos crudos}

\FunctionTok{print}\NormalTok{(}\StringTok{"Estructura inicial de los datos crudos:"}\NormalTok{)}
\end{Highlighting}
\end{Shaded}

\begin{verbatim}
## [1] "Estructura inicial de los datos crudos:"
\end{verbatim}

\begin{Shaded}
\begin{Highlighting}[]
\FunctionTok{str}\NormalTok{(raw\_data\_alcohol) }\CommentTok{\# Verificación inicial de la estructura y tipos de variables}
\end{Highlighting}
\end{Shaded}

\begin{verbatim}
## 'data.frame':    4324 obs. of  9 variables:
##  $ country                  : chr  "Italy" "Tanzania, United Republic of" "Cyprus" "Sierra Leone" ...
##  $ iso3                     : chr  "ITA" "TZA" "CYP" "SLE" ...
##  $ year                     : int  2003 2006 2017 2019 2017 2007 2018 2005 2000 2009 ...
##  $ sex                      : chr  "Both sexes" "Both sexes" "Both sexes" "Both sexes" ...
##  $ alcohol_liters_per_capita: num  9.623 6.338 6.117 0.253 6.024 ...
##  $ lower_ci                 : num  7.965 4.768 4.627 0.101 4.682 ...
##  $ upper_ci                 : num  11.389 8.109 7.798 0.567 7.537 ...
##  $ ci_width                 : num  3.424 3.342 3.171 0.467 2.855 ...
##  $ IndicatorCode            : chr  "SA_0000001688" "SA_0000001688" "SA_0000001688" "SA_0000001688" ...
\end{verbatim}

\begin{Shaded}
\begin{Highlighting}[]
\CommentTok{\#Procesamiento Encadenado}
\NormalTok{clean\_data\_alcohol }\OtherTok{\textless{}{-}}\NormalTok{ raw\_data\_alcohol }\SpecialCharTok{\%\textgreater{}\%}
 
  \FunctionTok{filter}\NormalTok{(sex }\SpecialCharTok{==} \StringTok{"Both sexes"} \SpecialCharTok{\&}\NormalTok{ year }\SpecialCharTok{\textgreater{}=} \DecValTok{2015} \SpecialCharTok{\&}\NormalTok{ year }\SpecialCharTok{\textless{}=} \DecValTok{2019}\NormalTok{) }\SpecialCharTok{\%\textgreater{}\%}  \CommentTok{\# Seleccionamos exclusivamente "Both sexes" y la ventana 2015{-}2019}
  
  \FunctionTok{select}\NormalTok{(country, iso3, year, alcohol\_liters\_per\_capita, }
\NormalTok{         lower\_ci, upper\_ci, ci\_width) }\SpecialCharTok{\%\textgreater{}\%} \CommentTok{\# Conservamos solo las variables necesarias para el estudio }
  
  \CommentTok{\# Gestión de valores faltantes (NA) }
  \FunctionTok{filter}\NormalTok{(}\SpecialCharTok{!}\FunctionTok{is.na}\NormalTok{(alcohol\_liters\_per\_capita) }\SpecialCharTok{\&} \SpecialCharTok{!}\FunctionTok{is.na}\NormalTok{(ci\_width)) }\SpecialCharTok{\%\textgreater{}\%}
  
  \CommentTok{\# Recodificación para asignar 5 Macro{-}Regiones \# Incluimos los códigos ISO3 de todos estos países}
  \FunctionTok{mutate}\NormalTok{(}\AttributeTok{Region =} \FunctionTok{case\_when}\NormalTok{(}
    
    \CommentTok{\# América Latina y el Caribe (Población hispanohablante + Brasil)}
\NormalTok{    iso3 }\SpecialCharTok{\%in\%} \FunctionTok{c}\NormalTok{(}\StringTok{"ATG"}\NormalTok{, }\StringTok{"ARG"}\NormalTok{, }\StringTok{"BHS"}\NormalTok{, }\StringTok{"BRB"}\NormalTok{, }\StringTok{"BLZ"}\NormalTok{, }\StringTok{"BOL"}\NormalTok{, }\StringTok{"BRA"}\NormalTok{, }\StringTok{"CHL"}\NormalTok{, }\StringTok{"COL"}\NormalTok{, }\StringTok{"CRI"}\NormalTok{, }\StringTok{"CUB"}\NormalTok{, }\StringTok{"DMA"}\NormalTok{, }\StringTok{"DOM"}\NormalTok{, }\StringTok{"ECU"}\NormalTok{, }\StringTok{"SLV"}\NormalTok{, }\StringTok{"GRD"}\NormalTok{, }\StringTok{"GTM"}\NormalTok{, }\StringTok{"GUY"}\NormalTok{, }\StringTok{"HTI"}\NormalTok{, }\StringTok{"HND"}\NormalTok{, }\StringTok{"JAM"}\NormalTok{, }\StringTok{"MEX"}\NormalTok{, }\StringTok{"NIC"}\NormalTok{, }\StringTok{"PAN"}\NormalTok{, }\StringTok{"PRY"}\NormalTok{, }\StringTok{"PER"}\NormalTok{, }\StringTok{"KNA"}\NormalTok{, }\StringTok{"LCA"}\NormalTok{, }\StringTok{"VCT"}\NormalTok{, }\StringTok{"SUR"}\NormalTok{, }\StringTok{"TTO"}\NormalTok{, }\StringTok{"URY"}\NormalTok{, }\StringTok{"VEN"}\NormalTok{) }\SpecialCharTok{\textasciitilde{}} \StringTok{"América Latina y el Caribe"}\NormalTok{,}
    
    \CommentTok{\# América del Norte (Solo Canadá y USA, ya que México está en ALC)}
\NormalTok{    iso3 }\SpecialCharTok{\%in\%} \FunctionTok{c}\NormalTok{(}\StringTok{"CAN"}\NormalTok{, }\StringTok{"USA"}\NormalTok{) }\SpecialCharTok{\textasciitilde{}} \StringTok{"América del Norte"}\NormalTok{,}
    
    \CommentTok{\# Asia (Sureste Asiático + Pacífico Occidental)}
\NormalTok{    iso3 }\SpecialCharTok{\%in\%} \FunctionTok{c}\NormalTok{(}\StringTok{"AFG"}\NormalTok{, }\StringTok{"BHR"}\NormalTok{, }\StringTok{"BGD"}\NormalTok{, }\StringTok{"BTN"}\NormalTok{, }\StringTok{"BRN"}\NormalTok{, }\StringTok{"KHM"}\NormalTok{, }\StringTok{"CHN"}\NormalTok{, }\StringTok{"COK"}\NormalTok{, }\StringTok{"CYP"}\NormalTok{, }\StringTok{"PRK"}\NormalTok{, }\StringTok{"EGY"}\NormalTok{, }\StringTok{"FJI"}\NormalTok{, }\StringTok{"IND"}\NormalTok{, }\StringTok{"IDN"}\NormalTok{, }\StringTok{"IRN"}\NormalTok{, }\StringTok{"IRQ"}\NormalTok{, }\StringTok{"ISR"}\NormalTok{, }\StringTok{"JPN"}\NormalTok{, }\StringTok{"JOR"}\NormalTok{, }\StringTok{"KAZ"}\NormalTok{, }\StringTok{"KGZ"}\NormalTok{, }\StringTok{"KIR"}\NormalTok{, }\StringTok{"KWT"}\NormalTok{, }\StringTok{"LAO"}\NormalTok{, }\StringTok{"LBN"}\NormalTok{, }\StringTok{"MYS"}\NormalTok{, }\StringTok{"MDV"}\NormalTok{, }\StringTok{"MHL"}\NormalTok{, }\StringTok{"FSM"}\NormalTok{, }\StringTok{"MNG"}\NormalTok{, }\StringTok{"MMR"}\NormalTok{, }\StringTok{"NRU"}\NormalTok{, }\StringTok{"NPL"}\NormalTok{, }\StringTok{"NZL"}\NormalTok{, }\StringTok{"NIU"}\NormalTok{, }\StringTok{"OMN"}\NormalTok{, }\StringTok{"PAK"}\NormalTok{, }\StringTok{"PLW"}\NormalTok{, }\StringTok{"PNG"}\NormalTok{, }\StringTok{"PHL"}\NormalTok{, }\StringTok{"QAT"}\NormalTok{, }\StringTok{"KOR"}\NormalTok{, }\StringTok{"WSM"}\NormalTok{, }\StringTok{"SAU"}\NormalTok{, }\StringTok{"SGP"}\NormalTok{, }\StringTok{"SLB"}\NormalTok{, }\StringTok{"LKA"}\NormalTok{, }\StringTok{"SYR"}\NormalTok{, }\StringTok{"TJK"}\NormalTok{, }\StringTok{"THA"}\NormalTok{, }\StringTok{"TLS"}\NormalTok{, }\StringTok{"TON"}\NormalTok{, }\StringTok{"TUR"}\NormalTok{, }\StringTok{"TKM"}\NormalTok{, }\StringTok{"TUV"}\NormalTok{, }\StringTok{"ARE"}\NormalTok{, }\StringTok{"UZB"}\NormalTok{, }\StringTok{"VUT"}\NormalTok{, }\StringTok{"VNM"}\NormalTok{, }\StringTok{"YEM"}\NormalTok{) }\SpecialCharTok{\textasciitilde{}} \StringTok{"Asia"}\NormalTok{,}
    
    \CommentTok{\# Europa }
\NormalTok{    iso3 }\SpecialCharTok{\%in\%} \FunctionTok{c}\NormalTok{(}\StringTok{"ALB"}\NormalTok{, }\StringTok{"AND"}\NormalTok{, }\StringTok{"ARM"}\NormalTok{, }\StringTok{"AUT"}\NormalTok{, }\StringTok{"AZE"}\NormalTok{, }\StringTok{"BLR"}\NormalTok{, }\StringTok{"BEL"}\NormalTok{, }\StringTok{"BIH"}\NormalTok{, }\StringTok{"BGR"}\NormalTok{, }\StringTok{"HRV"}\NormalTok{, }\StringTok{"CZE"}\NormalTok{, }\StringTok{"DNK"}\NormalTok{, }\StringTok{"EST"}\NormalTok{, }\StringTok{"FIN"}\NormalTok{, }\StringTok{"FRA"}\NormalTok{, }\StringTok{"GEO"}\NormalTok{, }\StringTok{"DEU"}\NormalTok{, }\StringTok{"GRC"}\NormalTok{, }\StringTok{"HUN"}\NormalTok{, }\StringTok{"ISL"}\NormalTok{, }\StringTok{"IRL"}\NormalTok{, }\StringTok{"ISR"}\NormalTok{, }\StringTok{"ITA"}\NormalTok{, }\StringTok{"KAZ"}\NormalTok{, }\StringTok{"KGZ"}\NormalTok{, }\StringTok{"LVA"}\NormalTok{, }\StringTok{"LTU"}\NormalTok{, }\StringTok{"LUX"}\NormalTok{, }\StringTok{"MLT"}\NormalTok{, }\StringTok{"MCO"}\NormalTok{, }\StringTok{"MNE"}\NormalTok{, }\StringTok{"NLD"}\NormalTok{, }\StringTok{"MKD"}\NormalTok{, }\StringTok{"NOR"}\NormalTok{, }\StringTok{"POL"}\NormalTok{, }\StringTok{"PRT"}\NormalTok{, }\StringTok{"MDA"}\NormalTok{, }\StringTok{"ROU"}\NormalTok{, }\StringTok{"RUS"}\NormalTok{, }\StringTok{"SMR"}\NormalTok{, }\StringTok{"SRB"}\NormalTok{, }\StringTok{"SVK"}\NormalTok{, }\StringTok{"SVN"}\NormalTok{, }\StringTok{"ESP"}\NormalTok{, }\StringTok{"SWE"}\NormalTok{, }\StringTok{"CHE"}\NormalTok{, }\StringTok{"TJK"}\NormalTok{, }\StringTok{"TUR"}\NormalTok{, }\StringTok{"TKM"}\NormalTok{, }\StringTok{"UKR"}\NormalTok{, }\StringTok{"GBR"}\NormalTok{, }\StringTok{"UZB"}\NormalTok{) }\SpecialCharTok{\textasciitilde{}} \StringTok{"Europa"}\NormalTok{,}
    
    \CommentTok{\# África }
\NormalTok{    iso3 }\SpecialCharTok{\%in\%} \FunctionTok{c}\NormalTok{(}\StringTok{"DZA"}\NormalTok{, }\StringTok{"AGO"}\NormalTok{, }\StringTok{"BEN"}\NormalTok{, }\StringTok{"BWA"}\NormalTok{, }\StringTok{"BFA"}\NormalTok{, }\StringTok{"BDI"}\NormalTok{, }\StringTok{"CMR"}\NormalTok{, }\StringTok{"CPV"}\NormalTok{, }\StringTok{"CAF"}\NormalTok{, }\StringTok{"TCD"}\NormalTok{, }\StringTok{"COM"}\NormalTok{, }\StringTok{"COG"}\NormalTok{, }\StringTok{"CIV"}\NormalTok{, }\StringTok{"COD"}\NormalTok{, }\StringTok{"DJI"}\NormalTok{, }\StringTok{"EGY"}\NormalTok{, }\StringTok{"GNQ"}\NormalTok{, }\StringTok{"ERI"}\NormalTok{, }\StringTok{"ETH"}\NormalTok{, }\StringTok{"GAB"}\NormalTok{, }\StringTok{"GMB"}\NormalTok{, }\StringTok{"GHA"}\NormalTok{, }\StringTok{"GIN"}\NormalTok{, }\StringTok{"GNB"}\NormalTok{, }\StringTok{"KEN"}\NormalTok{, }\StringTok{"LSO"}\NormalTok{, }\StringTok{"LBR"}\NormalTok{, }\StringTok{"LBY"}\NormalTok{, }\StringTok{"MDG"}\NormalTok{, }\StringTok{"MWI"}\NormalTok{, }\StringTok{"MLI"}\NormalTok{, }\StringTok{"MRT"}\NormalTok{, }\StringTok{"MUS"}\NormalTok{, }\StringTok{"MAR"}\NormalTok{, }\StringTok{"MOZ"}\NormalTok{, }\StringTok{"NAM"}\NormalTok{, }\StringTok{"NER"}\NormalTok{, }\StringTok{"NGA"}\NormalTok{, }\StringTok{"RWA"}\NormalTok{, }\StringTok{"STP"}\NormalTok{, }\StringTok{"SEN"}\NormalTok{, }\StringTok{"SYC"}\NormalTok{, }\StringTok{"SLE"}\NormalTok{, }\StringTok{"SOM"}\NormalTok{, }\StringTok{"ZAF"}\NormalTok{, }\StringTok{"SSD"}\NormalTok{, }\StringTok{"SDN"}\NormalTok{, }\StringTok{"SYR"}\NormalTok{, }\StringTok{"TGO"}\NormalTok{, }\StringTok{"TUN"}\NormalTok{, }\StringTok{"UGA"}\NormalTok{, }\StringTok{"TZA"}\NormalTok{, }\StringTok{"ZMB"}\NormalTok{, }\StringTok{"ZWE"}\NormalTok{) }\SpecialCharTok{\textasciitilde{}} \StringTok{"África"}\NormalTok{,}
    
    \CommentTok{\# Caso por defecto para códigos no reconocidos o Oceanía pura}
    \ConstantTok{TRUE} \SpecialCharTok{\textasciitilde{}} \StringTok{"Otras/No Clasificadas"}
\NormalTok{  )) }\SpecialCharTok{\%\textgreater{}\%}
  
  \CommentTok{\# Convertimos variables categóricas a factores para análisis grupal}
  \CommentTok{\# year lo convertimos a factor para comparación temporal}
  \FunctionTok{mutate}\NormalTok{(}
    \AttributeTok{country =} \FunctionTok{as.factor}\NormalTok{(country),}
    \AttributeTok{iso3 =} \FunctionTok{as.factor}\NormalTok{(iso3),}
    \AttributeTok{year =} \FunctionTok{as.factor}\NormalTok{(year), }\CommentTok{\# Convertimos año a factor para comparación temporal}
    \AttributeTok{Region =} \FunctionTok{as.factor}\NormalTok{(Region)}
\NormalTok{  )}

\CommentTok{\# Imprimimos las dimensiones del data frame limpio y las primeras filas}
\FunctionTok{print}\NormalTok{(}\FunctionTok{paste}\NormalTok{(}\StringTok{"Registros válidos (Ambos sexos, 2015{-}2019):"}\NormalTok{, }\FunctionTok{nrow}\NormalTok{(clean\_data\_alcohol)))}
\end{Highlighting}
\end{Shaded}

\begin{verbatim}
## [1] "Registros válidos (Ambos sexos, 2015-2019): 940"
\end{verbatim}

\begin{Shaded}
\begin{Highlighting}[]
\FunctionTok{print}\NormalTok{(}\StringTok{"Estructura final de los datos procesados y estratificados (str):"}\NormalTok{)}
\end{Highlighting}
\end{Shaded}

\begin{verbatim}
## [1] "Estructura final de los datos procesados y estratificados (str):"
\end{verbatim}

\begin{Shaded}
\begin{Highlighting}[]
\FunctionTok{str}\NormalTok{(clean\_data\_alcohol)}
\end{Highlighting}
\end{Shaded}

\begin{verbatim}
## 'data.frame':    940 obs. of  8 variables:
##  $ country                  : Factor w/ 188 levels "Afghanistan",..: 45 150 135 130 31 152 108 173 179 140 ...
##  $ iso3                     : Factor w/ 188 levels "AFG","AGO","ALB",..: 43 150 136 133 35 157 118 172 62 144 ...
##  $ year                     : Factor w/ 5 levels "2015","2016",..: 3 5 3 4 4 2 4 4 4 5 ...
##  $ alcohol_liters_per_capita: num  6.117 0.253 6.024 0.104 7.195 ...
##  $ lower_ci                 : num  4.6272 0.1009 4.6817 0.0364 5.3801 ...
##  $ upper_ci                 : num  7.798 0.567 7.537 0.34 9.066 ...
##  $ ci_width                 : num  3.171 0.467 2.855 0.304 3.686 ...
##  $ Region                   : Factor w/ 6 levels "África","América del Norte",..: 4 1 4 4 1 5 1 4 5 5 ...
\end{verbatim}

\begin{Shaded}
\begin{Highlighting}[]
\FunctionTok{head}\NormalTok{(clean\_data\_alcohol }\SpecialCharTok{\%\textgreater{}\%} \FunctionTok{select}\NormalTok{(country, iso3, year, Region))}
\end{Highlighting}
\end{Shaded}

\begin{verbatim}
##        country iso3 year Region
## 1       Cyprus  CYP 2017   Asia
## 2 Sierra Leone  SLE 2019 África
## 3  Philippines  PHL 2017   Asia
## 4     Pakistan  PAK 2018   Asia
## 5     Cameroon  CMR 2018 África
## 6     Slovakia  SVK 2016 Europa
\end{verbatim}

\subsection{Item 2}

Descripción\\

\subsubsection{Item 2.2}

Descripción\\

\subsection{Item n}

Descripción\\

\chapter{Análisis de Resultados}

\{\small Este capítulo presenta los resultados obtenidos al aplicar las
metodologías descritas en el capítulo anterior, así como del análisis
estadístico descriptivo de las variables estudiadas y las ventajas y
desventajas al aplicar dichos tests.\}

\section{Presentación de resultados 1}

Descripción\\

\section{Presentación de resultados 2}

Descripción\\

\section{Presentación de resultados n}

Descripción\\

\chapter{Conclusiones y Recomendaciones}

\textbf{1.-} Conclusion 1 del trabajo de investigación.\\

\textbf{Recomendación} Recomendación 1 del trabajo de investigación.\\

\textbf{2.-} Conclusion 2 del trabajo de investigación.\\

\textbf{Recomendación} Recomendación 2 del trabajo de investigación.\\

\textbf{n.-} Conclusion n del trabajo de investigación.\\

\textbf{Recomendación} Recomendación n del trabajo de investigación.\\

\appendix
\renewcommand{\thechapter}{\Alph{chapter}}

\chapter{Anexo A}

\chapter{Anexo B}

\chapter{Anexo C}

%%%%%%%%%%%%%%%%%%%%%%%%%%%%%%%%%%%%%%%%%%%%%%%%%%%%%
%                   REFERENCIAS                     %
%%%%%%%%%%%%%%%%%%%%%%%%%%%%%%%%%%%%%%%%%%%%%%%%%%%%%
% Mantenemos la estructura natbib de tu plantilla
\renewcommand{\bibname}{Lista de Referencias}
\bibliographystyle{apalike} 
\bibliography{bibliografia.bib} 

%%%%%%%%%%%%%%%%%%%%%%%%%%%%%%%%%%%%%%%%%%%%%%%%%%%%%
%                   APÉNDICES                       %
%%%%%%%%%%%%%%%%%%%%%%%%%%%%%%%%%%%%%%%%%%%%%%%%%%%%%

\end{document}
